\section{Statistical methods}
\label{app:statmethods}
\label{app:methods}

This appendix is the complete uncertainty-quantification methodology of the
measurement --- the training-seed variation, bootstrap and universe-covariance procedures,
the MAT-conformant covariance construction, the $1/\Phi$ flux normalization,
the comparison conventions ($\chi^{2}$, pulls, shape-only), the
ill-conditioning/rank analysis, the tension anatomy, the coverage toys, and
the closure tests. Each procedure gives the recipe, the math, the
implementation pointer, and a worked example using a reported number. We
document only the analysis-specific numerical conventions, not a general
statistics primer; the $\chi^{2}$ convention referenced throughout is
Eq.~\eqref{eq:chi2} in \S\ref{sec:chisq}.
The variance source this note calls \emph{training-seed variation} is named \texttt{seedscan} in the scripts, figure files and stored artifacts; the code name is retained there as provenance and the two refer to the same object.


\subsection{Numerical conventions: inverting a covariance (Cholesky, ridge jitter, pseudo-inverse)}
\label{sec:invert}

Equation~\eqref{eq:chi2} requires inverting $\mathbf{C}$. Three
practical issues come up in this analysis.

\paragraph{Positive-definiteness check (Cholesky).}
A matrix $\mathbf{C}$ is positive-definite (PD) iff there is a
lower-triangular $\mathbf{L}$ with $\mathbf{C}=\mathbf{L}\mathbf{L}\T$
(the Cholesky factorization). If Cholesky succeeds, $\mathbf{C}$ is
invertible and Eq.~\eqref{eq:chi2} is well-posed.

\paragraph{Ridge jitter.}
Numerical noise can leave a PD matrix barely non-PD. The standard fix
is to add an isotropic ridge $\varepsilon\, \mathbf{I}$. The
\texttt{uq/analyze\_uq.py} script does exactly this:
\begin{equation}
\varepsilon \;=\; 10^{-12} \,\frac{\Tr\mathbf{C}}{n},
\end{equation}
where the trace sets the natural scale of the eigenvalues. Twelve
orders of magnitude below the mean diagonal is far below any real
signal in $\mathbf{C}$ and safely above floating-point round-off.

\paragraph{Rank deficiency at $N \le n$.}
If you estimate $\mathbf{C}$ from $N$ replicas in a space of $n$ bins
with $N\le n$, the resulting sample covariance has rank at most
$N-1$: there are not enough independent ``directions'' to fill the
full $n\times n$ matrix. This was the regime of the Stage-1 bootstrap
($N=50$ replicas, $n=205$ bins): Cholesky failed without jitter, and
the required jitter, $1.6\times10^{-94}$, was a numerical
artifact of the rank deficiency rather than a real near-singularity.
The Stage-2 bootstrap ($N=300 > 205$) is full-rank. The
\emph{systematic} covariance suffers a more stubborn version of this
problem that survives at any universe count --- see \S\ref{sec:rank}.

\paragraph{Pseudo-inverse.}
For the $\chi^{2}$ comparison against the paper covariance,
\texttt{compare\_to\_paper\_fullcov.py} replaces $\mathbf{C}^{-1}$
with the Moore--Penrose pseudo-inverse $\mathbf{C}^{+}$, computed
via SVD. Writing $\mathbf{C}=\mathbf{U}\Sigma\mathbf{V}\T$ with
$\Sigma=\diag(\sigma_{1},\dots,\sigma_{n})$,
\begin{equation}
\mathbf{C}^{+} \;=\; \mathbf{V}\,\Sigma^{+}\,\mathbf{U}\T,
\qquad
\Sigma^{+}_{ii} \;=\;
\begin{cases}
1/\sigma_{i} & \sigma_{i} > \tau,\\
0 & \text{otherwise},
\end{cases}
\end{equation}
with a small threshold $\tau$ (NumPy's default is
$\tau = n \cdot \sigma_{\max} \cdot \varepsilon_{\text{machine}}$).
This silently drops singular modes that are numerically
indistinguishable from zero, replacing them with $0$ in the inverse
rather than $\infty$. The result is a stable quadratic form even when
$\mathbf{C}$ has zero or near-zero eigenvalues.

\subsection{Variance-source procedures}
\label{sec:variance}

Each subsection below follows the same template:
\emph{What it estimates $\cdot$ Procedure $\cdot$ Math $\cdot$ Code
pointer $\cdot$ Worked example}.

\subsubsection{ML stochasticity: training-seed variation}
\label{sec:seedscan}

\paragraph{What it estimates.}
The variance in the unfolded cross section due to random choices
inside the GBDT (gradient-boosted decision tree) optimizer. OmniFold
is deterministic given the input weights, the network architecture,
and the random seed; flipping only the seed produces slightly
different fits to the per-event reweighting functions, and that
spread is the ``ML noise floor''.

\paragraph{Procedure.}
Run $n=10$ full $5$-iteration MEFHC unfolds with everything fixed
except the GBDT \texttt{random\_state}. The driver
\texttt{sbatch\_unfold\_2d\_MEFHC\_5iter\_seedscan.sh} pins three
seeds per trial via \texttt{-{}-seed N}:
\begin{equation}
(\text{step-1 classifier},\text{step-2 classifier},\text{miss regressor})
\;=\;
(N,\,N+1,\,N+2),
\qquad N \in \{1,\dots,10\}.
\end{equation}
The canonical training-seed variation uses the \texttt{lgbm} (LightGBM) backend, so
it is backend-matched to the production unfold.

\paragraph{Math.}
For each reported bin $i$, build the $n$-vector $X_{i}^{(1)},\dots,
X_{i}^{(n)}$ across trials and compute the unbiased sample standard
deviation
\begin{equation}
\sigma^{\ML}_{i}
\;=\;
\sqrt{\frac{1}{n-1}\sum_{k=1}^{n}\big(X^{(k)}_{i}-\bar X_{i}\big)^{2}},
\qquad
\bar X_{i} = \frac{1}{n}\sum_{k} X^{(k)}_{i}.
\end{equation}
The headline diagnostic is the per-bin \emph{relative} spread
$\sigma^{\ML}_{i}/\bar X_{i}$, summarized by its median over the
$205$ reported bins.

\paragraph{Code pointer.}
\texttt{seedscan/analyze\_seedscan.py} (rollup);
\texttt{sbatch\_unfold\_2d\_MEFHC\_5iter\_seedscan.sh} (driver);
covariance output \texttt{uq/seedscan\_lgbm\_ml/uq\_covariance\_ml.root}.

\paragraph{Worked example (UQ envelope).}
With $n=10$ \texttt{lgbm} trials on the $205$-bin set:
\begin{itemize}[nosep]
  \item Total-$\sigma$ relative spread: $\boxed{0.0055\%}$.
  \item Per-bin median relative spread: $\boxed{0.166\%}$.
  \item Component trace $\sqrt{\Tr\mathbf{C}^{\ML}} = 5.061\times10^{-41}$.
\end{itemize}
The paper's per-bin total-uncertainty median is $6.86\%$; the ML
noise is therefore $\approx 2.4\%$ of the paper's quoted uncertainty
and far subdominant to data-stat and systematics. Going from $n=4$ to
$n=10$ trials moved the headline numbers within rounding, so the
training-seed variation envelope is converged.

\subsubsection{Statistical uncertainty: the Poisson bootstrap}
\label{sec:bootstrap}

\paragraph{What it estimates.}
The combined data-statistical and MC-statistical variance, propagated
\emph{coherently} through the full OmniFold pipeline. Coherent means:
if some data event happens to be ``up-weighted'' in a replica, the
same event flows through both the step-1 classifier training and the
final reweighted-MC cross-section construction; if an MC event is
up-weighted, the same up-weight is applied to its reco-tree and
truth-tree copies, preserving the response matrix.

\paragraph{Procedure.}
For each replica $k=1,\dots,N$:
\begin{enumerate}[nosep]
  \item Draw per-event Poisson(1) factors $b_{i}^{\,\text{data},(k)}$
        for every data event and $b_{i}^{\,\text{MC},(k)}$ for every
        MC event, using two independent sub-RNGs (a data RNG seeded
        from $k$, and an MC RNG seeded from $k + 10^{7}$). The two
        sources are uncorrelated by construction.
  \item Multiply data weights by $b^{\,\text{data}}$ and MC weights
        (both reco and truth copies of each MC event) by the
        \emph{same} $b^{\,\text{MC}}$. This keeps the migration
        matrix self-consistent.
  \item Run OmniFold (5 iterations) with the GBDT seed \emph{pinned}.
        Pinning the seed isolates the bootstrap signal from the ML
        noise of \S\ref{sec:seedscan}.
  \item Save \texttt{hXSec2D} (and 1D projections).
\end{enumerate}
Stage-1 used $N=50$ replicas (rank-deficient on $205$ bins); Stage-2
scaled to $N=300 > 205$ for a full-rank statistical covariance.

\paragraph{Caveat (seed-varying replicas).}
The $N=300$ covariance was produced with replicas that
varied \emph{both} \texttt{-{}-bootstrap-seed N} and \texttt{-{}-seed N},
so it carries a small ML-stochastic contribution on top of the pure
Poisson statistics. The
leak is numerically small relative to the systematic block (the ML
floor is $0.166\%$ vs the stat $0.564\%$); a purely Poisson block would
fix \texttt{-{}-seed} while varying only \texttt{-{}-bootstrap-seed}.

\paragraph{Math.}
Flatten each replica's $14\times16$ histogram into a $205$-vector
$\mathbf{X}^{(k)}$ over reported bins (row-major over $(\pT,\pPar)$).
The bootstrap covariance is the unbiased sample covariance:
\begin{equation}
\big(\mathbf{C}^{\boot}\big)_{ij}
\;=\;
\frac{1}{N-1}\sum_{k=1}^{N}
\big(X_{i}^{(k)} - \bar X_{i}\big)\big(X_{j}^{(k)} - \bar X_{j}\big),
\qquad
\bar X_{i} = \frac{1}{N}\sum_{k} X_{i}^{(k)}.
\label{eq:bootcov}
\end{equation}
This is exactly \texttt{np.cov(Xrep, rowvar=False, ddof=1)} in
\texttt{uq/analyze\_uq.py}. The rank-deficiency check is the
Cholesky-plus-jitter pattern from \S\ref{sec:invert}.

\paragraph{Why Poisson(1) is the right primitive.}
A non-parametric bootstrap with replacement assigns event $i$ a
multinomial count $n_{i}$ with $\mathbb{E}[n_{i}]=1$,
$\mathrm{Var}[n_{i}] = (N-1)/N \to 1$, and $\mathbb{E}[\sum n_{i}] =
N$. In the large-$N$ limit the $n_{i}$ become independent and each is
Poisson(1) (Appendix~\ref{app:poisson}). For a weighted analysis you
multiply $w_{i}\to w_{i}n_{i}$; since $n_{i}$ may equal zero,
some events are dropped from a given replica, exactly mimicking the
``with replacement'' resample.

\paragraph{Code pointer.}
Bootstrap weight draws and seed pinning:
\texttt{unfold\_2d\_omnifold\_unbinned.py} (\texttt{-{}-bootstrap-seed},
\texttt{-{}-seed}).\\
Covariance rollup: \texttt{uq/analyze\_uq.py}.\\
Output: \texttt{uq/bootstrap\_MEFHC\_300/uq\_covariance\_boot300.root}.

\paragraph{Worked example ($N=50$ vs $N=300$, MEFHC lgbm).}
\begin{center}
\begin{tabular}{@{}lcc@{}}
\toprule
Metric (205 bins) & $N=50$ & $N=300$ \\
\midrule
Total $\sigma$ relative std & $0.068\%$ & $\boxed{0.061\%}$ \\
Per-bin median relative spread & $0.532\%$ & $\boxed{0.564\%}$ \\
Per-bin $p_{84}$ & $1.235\%$ & $1.215\%$ \\
\bottomrule
\end{tabular}
\end{center}
The per-bin median is stable from $N=50$ to $N=300$ ($0.53\%\to
0.56\%$); the $N=300$ block is the one used in the headline because it
is full rank ($N>205$), whereas $N=50$ was rank-deficient and its
Cholesky needed jitter $1.6\times 10^{-94}$. The bootstrap median is
$\approx 3.4\times$ the training-seed variation median ($0.166\%$), as expected for a
sample of $\mathcal{O}(10^{6})$ data events. Across datasets the
per-bin median scales as $\sqrt{N_{\text{events}}}$: playlist $1A$
($\approx N/12$ of MEFHC) gives $1.46\%$, consistent with
$\sqrt{12}\times 0.56\% \approx 1.9\%$ at the same backend.

\subsubsection{Systematic uncertainty: universes}
\label{sec:universes}

\paragraph{What it estimates.}
Model uncertainty, by re-running the full unfold under each of
several physically-motivated reweightings of the MC. Where the
bootstrap shakes the \emph{events}, the universes shake the
\emph{model}: a different flux $\to$ different per-event weights $\to$
different unfolded cross section.

\paragraph{Procedure.}
The C++ event loop (\texttt{runEventLoopOmniFold.cpp}) dumps, for each MC event,
alternative weight branches $w^{\text{truth}}_{B,u}$ and
$w^{\text{reco}}_{B,u}$ for every (band $B$, universe index $u$)
pair. The Stage-2 sweep has \textbf{187 universes across 44 bands}. Five bands
are genuinely kinematic: BeamAngleX/Y, MuonResolution, and
Muon\_Energy\_MINERvA/MINOS. They require selection-complete
\texttt{MNV101\_ACTIVE\_UNIVERSE} event loops when migration across the event
selection matters. MinosEfficiency and GEANT are weight-only. The
dominant ensemble is the $100$-universe PPFX flux band; almost all
other bands are $\pm1\sigma$ pairs.

\begin{center}
\begin{tabular}{@{}llr@{}}
\toprule
Band (examples) & Variation & $N_{u}$ \\
\midrule
\texttt{Flux} & PPFX hadron-production sampling & 100 \\
\texttt{Muon\_Energy\_MINOS} & Muon energy scale, $\pm 1\sigma$ & 2 \\
\texttt{Muon\_Energy\_MINERvA} & Muon energy scale (lateral), $\pm1\sigma$ & 2 \\
\texttt{MinosEfficiency} & MINOS muon-match efficiency, $\pm 1\sigma$ & 2 \\
\texttt{MaRES},\,\texttt{MvRES},\,\texttt{MaCCQE} & GENIE knobs, $\pm 1\sigma$ & 2 ea. \\
\texttt{GEANT\_\{Neutron,Pion,Proton\}} & FSI response (lateral), $\pm1\sigma$ & 2 ea. \\
\bottomrule
\end{tabular}
\end{center}
For each $(B,u)$, \texttt{unfold\_2d\_omnifold\_unbinned.py
-{}-universe B:u} reweights all MC events (truth and reco) by the
matching branch and runs OmniFold with a fixed seed. A separate
matched ``CV'' job runs without \texttt{-{}-universe}, with the
\emph{same} omnifile, backend, seed, iteration count, and flux, so
that ML noise cancels between CV and universe.

\paragraph{Math (MAT-conformant convention).}
For universe $(B,u)$, define the per-bin value
$\mathbf{X}^{(B,u)}$, a $205$-vector over reported bins. MINERvA's MAT
framework
(\texttt{PlotUtils/\allowbreak MnvVertErrorBand::\allowbreak CalcCovMx}, see \S\ref{sec:minerva})
builds the per-band covariance as the \emph{mean-centered, biased}
sample covariance, referenced to the \textbf{universe mean} (not the
CV) and normalized by $1/N_{u}$:
\begin{equation}
\boxed{\;
\mathbf{C}^{B}
\;=\;
\frac{1}{N_{u}}\sum_{u=1}^{N_{u}}
\big(\mathbf{X}^{(B,u)} - \langle\mathbf{X}\rangle_{B}\big)
\big(\mathbf{X}^{(B,u)} - \langle\mathbf{X}\rangle_{B}\big)\T,
\qquad
\langle\mathbf{X}\rangle_{B} = \frac{1}{N_{u}}\sum_{u}\mathbf{X}^{(B,u)}.
\;}
\label{eq:matcov}
\end{equation}
This single formula is applied \emph{uniformly} to all $N_{u}\ge 2$ ---
there is \textbf{no special-case $\pm1\sigma$ pair formula}. For a
$\pm1\sigma$ pair ($N_{u}=2$) it collapses to the rank-1
$\tfrac14(\mathbf{X}_{+}-\mathbf{X}_{-})(\mathbf{X}_{+}-\mathbf{X}_{-})\T$;
mean-centering kills the symmetric component, so a one-sided or
asymmetric pair correctly contributes near-zero variance rather than
the rank-2 contribution a naive CV-centered formula would give
(Appendix~\ref{app:pair}).

\paragraph{Target-nucleon normalization (rank-1 add-on).}
MINERvA quotes a $1.4\%$ overall target-nucleon normalization
uncertainty (Aliaga NIM, arXiv:1305.5199) that lives outside the
\texttt{MnvVertErrorBand} machinery. We add it as a fully-correlated
rank-1 outer product on top of the band sum:
\begin{equation}
\mathbf{C}^{\text{norm}} \;=\; (\sigma_{N}\,\mathbf{X}^{\CV})(\sigma_{N}\,\mathbf{X}^{\CV})\T,
\qquad \sigma_{N} = 0.014,
\label{eq:normband}
\end{equation}
via \texttt{analyze\_universes.py -{}-add-norm 0.014}.

\paragraph{Total systematic covariance.}
Bands are assumed independent (different physics sources, different
RNGs), so
\begin{equation}
\mathbf{C}^{\syst}
\;=\;
\sum_{B}\mathbf{C}^{B}\;+\;\mathbf{C}^{\text{norm}}.
\label{eq:systblocksum}
\end{equation}
Independence is a modeling assumption with one known correlated exception --- the
Bashyal flux\,$\leftrightarrow$\,Muon\_Energy\_MINOS joint block
(\S\ref{sec:rank}), which we have since rederived ($\rho=0.900$,
\texttt{uq/flux\_muonE\_cross.root}); it is real but rank-preserving, so for the
2D budget the block sum above is retained as the headline and the cross term is
carried as a diagnostic.

There is, however, a second and larger correction that Eq.~\eqref{eq:systblocksum}
omits: the band sum is linear in the per-band covariances, whereas the OmniFold
reweighter is \emph{nonlinear} in the input perturbations, so jointly perturbing
all bands and re-unfolding produces a cross term beyond any pairwise correlation.
We measure this directly with a \emph{unified throw} --- $N$ joint re-unfolds in
which every reweightable band and the flux multisim are perturbed simultaneously
per event ($w_{\CV}\prod_B\rho_B^{g_B}$, $g_B\sim\mathcal N(0,1)$) and the cross
section re-extracted.  The corrected implementation uses asymmetric $\pm1\sigma$
endpoint interpolation, one fixed estimator seed, and covariance about the throw
mean; the mean shift from CV is stored separately. It does not subtract a scalar
run-to-run jitter floor (\texttt{nd-unfolding/unified\_throw\_cov.py}). The 5D
result uses the strict joint ensemble and matched-endpoint comparator described
in \S\ref{sec:uq-combine}. No corrected 4D or extended-fiducial covariance is
reported. A single-throw superposition probe is insufficient to test the
nonlinear response. Any adoption must document the low-rank treatment, mean
shift, and selection-support bound before transferring information onto a sweep block,
$\mathbf{C}^{\syst}\!\to\!(\mathbf{C}^{\syst}-\mathbf{C}^{\mathrm{vert}})+\mathbf{G}\,\mathbf{C}^{\mathrm{vert}}\mathbf{G}$
($\mathbf{G}=\operatorname{diag}g$), which is positive-semidefinite by construction
(a direct rank-$N$ swap is not) and never falls below the block baseline.

\paragraph{Code pointer.}
Event-loop dump:
\texttt{runEventLoopOmniFold.cpp} (env \texttt{MNV101\_DUMP\_UNIVERSES}).\\
Per-universe unfold:
\texttt{unfold\_2d\_omnifold\_unbinned.py -{}-universe BAND:IDX}.\\
Per-band rollup: \texttt{uq/analyze\_universes.py} (\texttt{-{}-add-norm 0.014}).\\
Output:
\texttt{uq/universe\_stage2\_MEFHC\_full\_matcorr/uq\_universe\_covariance\_full\_matcorr.root}.

\paragraph{Worked example (full matcorr sweep).}
Top contributing bands, per-bin median relative
$\sigma_{i}/X^{\CV}_{i}$ over $205$ reported bins:
\begin{center}
\begin{tabular}{@{}lc@{}}
\toprule
Band & Per-bin median rel.\ $\sigma$ \\
\midrule
Muon\_Energy\_MINOS & $2.305\%$ \\
Muon\_Energy\_MINERvA & $1.285\%$ \\
MinosEfficiency & $1.472\%$ \\
Flux (PPFX, $N_{u}=100$) & $\dead{1.008\%}$ $\to$ $4.993\%$ \\
MaRES & $0.548\%$ \\
MvRES & $0.377\%$ \\
MaCCQE & $0.357\%$ \\
\midrule
Systematic total (44 bands $+$ norm) & $\dead{4.783\%} \to \boxed{6.830\%}$ (median); $\sqrt{\Tr\mathbf{C}}=\dead{2.463\times10^{-39}} \to 3.214\times10^{-39}$ \\
\bottomrule
\end{tabular}
\end{center}
\textbf{Struck entries are the pre-flux-fix rollup} (output path above,
\texttt{...\_full\_matcorr/}); the live values to their right are the flux-fixed
rollup \texttt{...\_full\_matcorr\_fluxfix/uq\_universe\_summary.txt}, which is the
product this note quotes (\S\ref{sec:systematics}). Of the individual systematic
bands only \texttt{Flux} moved; the grouped \emph{Statistical} category also moved
($0.563\% \to 0.549\%$) for the separate reason given in the provenance note in
\S\ref{sec:combine}.
Grouped by category the medians are: Muon reconstruction $3.366\%$,
Models $2.034\%$, Flux $\dead{1.008\%} \to 4.993\%$, Hadronic response $0.758\%$,
Statistical $\dead{0.563\%} \to 0.549\%$, Normalization $0.081\%$. Muon energy scale
(MINOS$+$MINERvA) and the MINOS-efficiency band dominate; these are
exactly the bands the Ruterbories paper flags as flux-correlated
(\S\ref{sec:rank}).

\subsubsection{How MINERvA builds its published covariance}
\label{sec:minerva}

MINERvA's covariance construction is not spelled out in the paper; it
is defined in the MAT C++
framework. The takeaway: \textbf{MAT samples per-band exactly as we do},
using the mean-centered $1/N$ sample covariance, and sums bands
element-wise.

\paragraph{What the Ruterbories 2021 paper says (Section V).}
Three uncertainty categories --- flux, detector response,
neutrino-interaction model --- evaluated by ``re-extracting the cross
section using modified simulations'', i.e.\ the universe method, but
with \emph{no covariance formula given}. Specifics quoted:
flux below $4\%$ across most of phase space thanks to the
$\nu$-$e$-scattering constraint (Valencia, arXiv:1906.00111);
$1.4\%$ target-nucleon normalization; in-situ test-beam hadron
response with per-particle GEANT inelastic variations; muon energy
scale $1\%$ constrained by the low-recoil fit of Bashyal
(JINST 16 P08068), with the explicit statement that
\emph{``the fit causes the flux and muon energy scale uncertainties to
be correlated and that correlation is propagated to the final
result''} --- the only non-block-diagonal piece the paper admits.
The only covariance that appears in equations is $V$ in the model-test
$\chi^{2}$ (Eq.~1, ``the measurement covariance matrix''), with $205$
degrees of freedom (Table I), and a log-normal $\chi^{2}$ variant
motivated by the highly-correlated flux normalization (Peelle's
Pertinent Puzzle).

\paragraph{What the paper does \emph{not} say.}
No per-band covariance construction; no $\pm1\sigma$ pair formula; no
band-independence/block-sum statement; nothing about rank, condition
number, or regularization of $V$; no bootstrap or coverage validation
(the paper uses D'Agostini iterative unfolding). All of this is
delegated to MAT or simply unstated.

\paragraph{Per-band covariance: \texttt{MnvVertErrorBand::CalcCovMx}.}
Reading the MAT source establishes three facts, encoded in
Eq.~\eqref{eq:matcov}:
\begin{itemize}[nosep]
  \item the reference is the \textbf{universe mean}
        $\langle x\rangle_{u}$, not the central value $X_{\CV}$;
  \item the normalization is \textbf{$1/N$} (biased), not Bessel's
        $1/(N-1)$ --- a $1\%$ offset for PPFX ($N=100$);
  \item there is \textbf{no special case for $N=2$} --- the
        mean-centered $N$-universe form is applied uniformly. The
        ``MINERvA-101 $\pm1\sigma$ pair'' formula
        $\tfrac12(\boldsymbol{\Delta}_{+}\boldsymbol{\Delta}_{+}\T +
        \boldsymbol{\Delta}_{-}\boldsymbol{\Delta}_{-}\T)$ is a folk
        re-derivation that agrees only for CV-symmetric shifts.
\end{itemize}
We verified our rollup against MAT byte-for-byte: populating a
\texttt{PlotUtils::MnvH1D} with the same $44$ bands / $188$ universes
and comparing \texttt{MnvH1D::GetTotalErrorMatrix} against our
\texttt{-{}-add-norm 0} rollup gives a max element-wise relative
difference of $5.5\times10^{-17}$ (machine $\varepsilon$), with
$\sqrt{\Tr\mathbf{C}}$ matching to seven digits
(\texttt{uq/verify\_matcorr\_vs\_mnvh1d.py}).

\paragraph{Total covariance: \texttt{MnvH1D::GetTotalErrorMatrix}.}
A straight element-wise sum across bands plus the statistical matrix.
Band-independence is assumed; cross-band correlations exist only if
encoded inside a single combined band (e.g.\ the Bashyal joint fit,
handled outside this loop). MAT does \emph{no} ridge / shrinkage / SVD
truncation on the published total cov and no special handling of rank
deficiency --- the published $V$ is full-rank because it includes the
\emph{statistical} block, not because it sums richer systematics (its
systematic-only rank is $120<140$; see \S\ref{sec:rank}).

\paragraph{Upstream references that set each band's $\pm1\sigma$ amplitude.}
\begin{center}
\begin{tabular}{@{}l@{\quad}p{0.56\textwidth}@{}}
\toprule
Input & Reference \\
\midrule
Flux / PPFX (100 universes) & Aliaga, PRD 94 092005 (2016), arXiv:1607.00704 \\
$\nu$-$e$ flux constraint ($<4\%$) & Valencia, PRD 100 092001 (2019), arXiv:1906.00111 \\
Muon-E scale $\leftrightarrow$ flux fit & Bashyal, JINST 16 (2021) P08068 \\
Test-beam / GEANT hadron response & Aliaga, NIM A 789 28 (2015), arXiv:1501.06431 \\
GENIE reweight knobs & Andreopoulos, NIM A 614 87 (2010); GENIE manual arXiv:1510.05494 \\
$1.4\%$ normalization (detector mass) & Aliaga, NIM A 743 130 (2014), arXiv:1305.5199 \\
\bottomrule
\end{tabular}
\end{center}

\subsubsection{Block-sum combination}
\label{sec:combine}

\paragraph{What it estimates.}
The total uncertainty covariance, combining the three statistically
independent components above.

\paragraph{Independence assumption.}
The bootstrap samples \emph{events} (with two distinct sub-RNGs for
data and MC). The universes vary \emph{model weights}, with the GBDT
seed pinned. The training-seed variation scans the \emph{GBDT seed}, with all
weights fixed. The three Monte Carlos are independent by
construction:
\begin{equation}
\mathbf{C}^{\tot}
\;=\;
\mathbf{C}^{\ML} \,+\, \mathbf{C}^{\stat} \,+\, \mathbf{C}^{\syst}.
\label{eq:blocksum}
\end{equation}
The combined-cov hook in \texttt{compare\_to\_paper\_fullcov.py}
(\texttt{-{}-omnifold-cov ROOT:HIST}, repeatable) adds OmniFold-
derived covariances on top of the paper's \texttt{TotalCovariance}
through exactly this rule. $\mathbf{C}^{\ML}$ is more than an order of
magnitude below the other two; the headline includes it but it is
numerically negligible.

\paragraph{Worked example (UQ envelope).}
Component traces and per-bin median relative spreads:
\begin{center}
\begin{tabular}{@{}lcc@{}}
\toprule
Component & $\sqrt{\Tr\mathbf{C}}$ & Per-bin median rel.\ $\sigma$ \\
\midrule
ML noise (lgbm training-seed variation, $n=10$) & $5.061\times10^{-41}$ & $0.166\%$ \\
Statistical (Poisson bootstrap, $N=300$) & $1.817\times10^{-40}$ & $0.549\%$ \\
Systematic (fluxfix sweep $+$ norm) & $3.214\times10^{-39}$ & $6.830\%$ \\
\midrule
\textbf{Combined (block sum)} & $3.220\times10^{-39}$ & $\boxed{6.865\%}$ \\
Paper \texttt{TotalCov} (reference) & $2.676\times10^{-39}$ & $6.86\%$ \\
\bottomrule
\end{tabular}
\end{center}
The combined OmniFold envelope essentially \emph{matches} the paper's per-bin
median uncertainty ($6.865\%$ against $6.86\%$), dominated by the systematic
block. (The combined $\sqrt{\Tr\mathbf{C}}$ \emph{is} the quadrature sum of the three
independent component traces. The per-bin median column is \textbf{not}: a median
of a sum is not the quadrature sum of the medians, so $6.865\%$ does not --- and
should not --- reproduce by combining the three rows above it. Both combined
entries are read from the block-sum table in
\texttt{2d-unfolding/2D\_OMNIFOLD\_STUDY\_STATUS.md}, \S``Stage-2 UQ envelope''.
The flux-fixed rollup summary itself reports the universe $+$ bootstrap
combination, before the ML block, as $6.845\%$
[\texttt{...\_matcorr\_fluxfix/uq\_universe\_summary.txt}, line
``\texttt{Combined (universe + bootstrap) median rel}''] --- also a read value,
not one derived from the entries printed here.)

\begin{quote}
\textbf{Provenance note (2026-07-31; table corrected 2026-08-11).} The table above
now carries the \emph{flux-fixed} rollup. What it previously printed, and what a
reader may remember from an earlier draft, was the \emph{pre-flux-fix} product ---
systematic $\dead{2.463\times10^{-39}}$ / $\dead{4.783\%}$, combined
$\dead{2.470\times10^{-39}}$ / $\dead{4.822\%}$ --- together with the inference drawn
from them, that the OmniFold envelope is $\dead{\approx\!70\%}$ of the paper's. Those
four figures and that inference are \textbf{retracted}: they are artefacts of a
superseded covariance, not of the measurement. The current budget agrees with the one
tabulated in \S\ref{sec:systematics}, where the envelope essentially \emph{matches}
the paper ($6.865\%$ against $6.86\%$).

One further cell of the table moved for an unrelated and earlier reason, recorded
here because it is a correction in its own right and \emph{not} because it changed
the headline: the statistical block is quoted from the pure-Poisson $N=300$
bootstrap ($\dead{1.828\times10^{-40}}$ / $\dead{0.564\%}$ $\to$
$1.817\times10^{-40}$ / $0.549\%$), which superseded a seed-varying set that leaked
ML stochasticity into the statistical block
(\texttt{2d-unfolding/2D\_OMNIFOLD\_STUDY\_STATUS.md}, 2026-05-29).
\textbf{It is not what moved the combined row.} The block sum rounds to
$3.220\times10^{-39}$ with either statistical cell --- the change sits about two
orders of magnitude below the printed precision --- and essentially the whole
movement of the combined entry comes from the systematic cell, i.e.\ from the flux
fix. Stated explicitly so that a later check does not reverse a correct edit on
finding that the headline is insensitive to it.

The two rollups differ in \emph{exactly one band}. Comparing
\texttt{2d-unfolding/uq/universe\_stage2\_MEFHC\_full\_matcorr/} against
\texttt{...\_matcorr\_fluxfix/}, every systematic band agrees to the last quoted
digit except \texttt{full\_Flux}, whose per-bin median rises
$1.008\% \rightarrow 4.993\%$ once each PPFX universe is divided by its own flux
integral $\Phi_u$ instead of by $\Phi_{\rm CV}$. That single correction carries the
universe total $4.783\% \rightarrow 6.830\%$ and the combined figure to
$6.845\%$ before the ML block, which is the whole of the discrepancy. The pull
widths later in this appendix (\S\ref{sec:pull}) were affected identically:
$\dead{0.069}/\dead{0.466}$ were computed on the same superseded covariance, and
the flux-fixed values now printed there are $0.051/0.409$
(\texttt{VALIDATION\_LEDGER.md}).

The same defect, uncorrected, is still present in the N-D and 5-D kernels; see
\texttt{AUDIT-FINDINGS-20260731.md} J18 and J28.
\end{quote}

\subsection{Comparing to the paper}
\label{sec:compare}

\subsubsection{The 205-bin selection}
\label{sec:205bins}

The paper reports $\mathrm{d}^{2}\sigma/(\mathrm{d}\pT\,\mathrm{d}\pPar)$
on a $14$-bin $\pT$ axis $\times$ $16$-bin $\pPar$ axis $=$ $224$
cells (edges in
\texttt{minerva\_paper\_anc/bin\_mapping.txt}). Of those, $205$ have
$x^{\paper}>0$ in
\texttt{data\_result\_ptpl\_2D\_minerva\_inclusive\_6GeV.txt}; the
remaining $19$ are paper-unreported cells that the paper's analysis
suppresses via an efficiency cut. We drop those $19$ from both sides
of every $\chi^{2}$ and pull --- that is the unambiguous
``apples-to-apples'' choice.

A secondary diagnostic mask, used only as a sensitivity check, is the
strict-interior set: cells where $\pT^{\text{high}} / \pPar^{\text{low}}
\le \tan 20^{\circ}$. Some cells in the $205$-bin set straddle the
$\theta_{\mu}<20^{\circ}$ kinematic cut and have wildly suppressed
paper efficiencies that we do not replicate; pulls on those cells run
to $\mathcal{O}(100)\sigma$ and dominate the $\chi^{2}$. The $185$-bin
interior cut removes them, and we quote it only as a sensitivity check that
the headline tension is not an artifact of those few boundary cells.
\textbf{The canonical metric is the 205-bin set} used everywhere below.

\subsubsection{Covariance-aware $\chi^{2}/\mathrm{ndf}$}
\label{sec:chisq}

\paragraph{Math.}
With $\boldsymbol{\Delta}\equiv\mathbf{x}^{\ours}-\mathbf{x}^{\paper}$
on the $205$-bin set and $\mathbf{C}$ the chosen covariance,
\begin{equation}
\label{eq:chi2}
\chi^{2}
\;=\;
\boldsymbol{\Delta}\T\,\mathbf{C}^{+}\,\boldsymbol{\Delta},
\qquad
\mathrm{ndf} \;=\; n_{\text{reported}} = 205,
\qquad
\chi^{2}/\mathrm{ndf} \;=\; \chi^{2}/205.
\end{equation}

\paragraph{This $\mathrm{ndf}$ is dimension-conditional, and is justified here by the truncation
scan --- not by convention.} Setting $\mathrm{ndf}=n_{\text{reported}}$ presumes the covariance has
close to $n_{\text{reported}}$ well-determined directions. On the $205$-bin 2D problem that
presumption is \emph{tested and holds}: the rank-truncation scan below rises smoothly
$0.69\,(r{=}50)\to2.35\,(100)\to3.30\,(180)\to3.66\,(205)$ with no plateau or late jump, the
signature of a spectrum whose effective rank is comparable to its dimension. The scan is the
evidence; $205$ is not assumed.

\textbf{It does not transfer to the N-D covariances.} The standard-P4 5D candidate covariance has
$10\,694$ reported bins and a measured effective positive rank of $\mathbf{263}$ (numerical
condition number infinite), and its $4825$-bin 5D$\to$4D projection preserves that rank exactly at
$263$. A sum of ${\sim}45$ two-endpoint MAT bands plus statistical and ML blocks cannot span
$10\,694$ directions, so this is a property of the construction, not a defect. Quoting
$\chi^{2}/n_{\text{reported}}$ against such a covariance divides by a number roughly $40\times$
(5D) or $18\times$ (4D) larger than the count of directions actually contributing to the quadratic
form, and would understate any tension by that factor.

\textbf{Therefore $\mathrm{ndf}=n_{\text{reported}}$ must not be used for an N-D covariance.} Any
N-D $\chi^{2}$ shall instead quote (i) the inverse actually used --- pseudo-inverse with its
\texttt{rcond}, or the truncation rank --- stated at the point of quotation rather than by
reference; (ii) the \emph{retained rank} as $\mathrm{ndf}$; (iii) the rank-truncation scan for that
covariance, which is expected \emph{not} to be smooth, and whose smoothness is precisely what
licensed the 2D choice above; and (iv) which covariance is meant, since the 5D candidate, its 4D
projection and the published 2D block have different ranks; and (v) \emph{for every
sample-covariance block entering the sum}, its ensemble size $N$, its normalization convention
(biased $1/N$ or unbiased $1/(N-1)$), the effective dimension $p$ actually inverted after
truncation, and the finite-ensemble treatment applied --- or the explicit statement that none was.
Measurements and derivation: \texttt{docs/orchestration/RANK-AND-INVERSION-20260810.md}; the
predeclared rank cut itself is a statistics decision recorded in
\texttt{docs/STATVAL\_REPAIR.md}.

\paragraph{Why declaration (v), and why it does not yet mandate a correction.}
Requirements (i)--(iv) can all be satisfied while a quoted $\mathrm{ndf}$ silently inherits an
uncorrected finite-ensemble bias from the sample-covariance blocks inside the sum, because not one
of them is the ensemble size of any such block. The inverse of a noisy sample covariance is biased
\emph{upward}: for an unbiased $1/(N-1)$ estimator over $N$ independent Gaussian realizations the
standard debiasing factor is $(N-p-2)/(N-1)$, and for this analysis's biased $1/N$ production
convention (\texttt{nd-unfolding/uq\_math.py}, ``universe-mean centered, biased 1/N'') the
corresponding factor is $(N-p-2)/N$. On the unified-throw 5D candidate, with $N=160$ joint throws
and a truncated $p=100$, that is $58/160\simeq0.363$ --- the uncorrected precision matrix is too
large by a factor $\simeq2.76$. \textbf{The consequence runs the opposite way to the intuition:} a
$\chi^{2}$ evaluated on a fixed residual is \emph{inflated}, not flattered, so tension is
overstated rather than hidden; what becomes over-optimistic is any confidence region derived from
the same over-tight covariance.
\textbf{No generic correction is mandated here, deliberately.} The quoted covariance is a
\emph{sum} of about 45 deterministic rank-one MAT band outer products plus statistical and ML
blocks, and a sum of that kind has no single debiasing factor; the standard factor also assumes
independent Gaussian realizations and a truncation dimension chosen independently of the data,
neither of which is established for a data-dependent rank cut. The treatment is therefore to be
decided \emph{for the covariance and inverse actually adopted}, once (v) makes its operands exist.
The cheap half is unconditional: record $N$ beside each block in the construction receipts.

\paragraph{What protects the \emph{quoted} combined $\chi^{2}$: an external statistical floor.}
The headline $\chiCombined$ does not invert our blocks alone. It inverts the paper$+$ours
\emph{sum}, and that sum carries the paper's \texttt{StatOnlyCov}, which is rank $205$ and
near-diagonal (\S\ref{sec:rank}); a full-rank, near-diagonal block puts a variance floor under
\emph{every} eigendirection of the combined matrix. It is also the larger of the two statistical
contributions, by a factor $\simeq2.5$ in \emph{square-root} trace and therefore by its square,
$2.55^{2}\simeq6.5$, in \emph{trace} (operands in \S\ref{sec:rank}). Naming the power matters:
the two figures describe the same comparison and differ by more than a factor of two. The
small-eigenvalue directions on which a finite-$N$ debiasing factor would act are consequently set
by an external, published block rather than by any ensemble of ours, and that is what bounds this
quoted number's exposure to declaration~(v).
\emph{This is deliberately not an argument from trace share.} A trace weights eigenvalues by
$\lambda$ while a precision matrix weights them by $1/\lambda$, so a block holding a negligible
fraction of the trace can still dominate an inverse. The ours-only construction is the
demonstration: remove the external floor and our finite-$N$ blocks become precisely what fills the
null space of $\mathbf{C}^{\syst}$ (rank $140\to201$), and the pseudo-inverse returns
$\chi^{2}/\mathrm{ndf}=252$ (\S\ref{sec:rank}). The protection therefore concerns \emph{where the
smallest retained eigenvalues come from}, not how much trace our blocks hold. It does not make
$\chiCombined$ a calibrated goodness-of-fit --- it is not one, for the separate shared-systematics
reason given with the worked example above.

The implementation (\texttt{chi2\_with\_cov} in
\texttt{compare\_to\_paper\_fullcov.py}) builds the $205\times205$
block, inverts with \texttt{np.linalg.pinv} (\S\ref{sec:invert}), and
forms the quadratic form.

\paragraph{Code pointer.}
\texttt{compare\_to\_paper\_fullcov.py} (headline);
\texttt{-{}-omnifold-cov ROOT:HIST} (repeatable) adds each OmniFold
covariance via Eq.~\eqref{eq:blocksum}; \texttt{-{}-subtract-stat}
removes the paper's stat block to avoid double-counting the bootstrap
(\S\ref{sec:rank}).

\paragraph{Worked example: comparison to paper (headline numbers).}
\begin{center}
\begin{tabular}{@{}lcc@{}}
\toprule
Metric (205 bins) & Paper & OmniFold (combined) \\
\midrule
Width-weighted total $\sigma$, $\sqrt{\mathbf{w}\T\mathbf{C}\,\mathbf{w}}/(\mathbf{w}\T\mathbf{x})$ & $4.61\%$ & $2.26\%$ \\
Per-bin median relative $\sigma$ & $6.86\%$ & $\dead{4.822\%} \to 6.865\%$ \\
$\chi^{2}/\mathrm{ndf}$ & $3.661$ (paper-only) & $\boxed{\chiCombined}$ (paper$+$ours) \\
Log-normal $\chi^{2}/\mathrm{ndf}$ & --- & $\chiCombinedLog$ \\
\bottomrule
\end{tabular}
\end{center}
The $3.661 \to \chiCombined$ drop ($-60\%$) is the central UQ
result: once the OmniFold-side error budget is included,
the summed-covariance distance to the paper falls to $\chi^{2}/\mathrm{ndf}
\approx 1.5$. (A perfectly correct model with correct uncertainties
gives exactly $1$; $\approx1.5$ tolerates small residual shape mismatch in
high-$\pPar$ tails, see \S\ref{sec:shape}.) \textbf{This is an indicative
distance, not a calibrated goodness-of-fit:} the published (IBU-era) result
and this OmniFold result share the same MINERvA data and flux/detector/GENIE
systematics, and without the (unavailable) OmniFold$\leftrightarrow$paper
cross-covariance the summed covariance both double-counts those shared
systematics and omits their correlation. We therefore quote it as a distance
and ratio-to-paper, not a $p$-value; see \S\ref{sec:rank} for why it remains
the headline and why we report per-bin pulls rather than a single unqualified
ours-only $\chi^{2}$.
The log-normal variant (matching the paper's Table I treatment of the
correlated flux normalization) gives $\chiCombinedLog$, essentially identical.

\paragraph{Edge-bin sensitivity ($\pPar$-min cut).}
\texttt{compare\_to\_paper\_fullcov.py} ships a side-table that
recomputes $\chi^{2}/\mathrm{ndf}$ on progressively shrinking subsets,
using the paper covariance only (no OmniFold add):
\begin{center}
\begin{tabular}{@{}rrrrr@{}}
\toprule
$\pPar \ge$ (GeV/c) & $N_{\text{bins}}$ & $\chi^{2}$ & $\chi^{2}/\mathrm{ndf}$ & median ratio \\
\midrule
1.5 & 205 & 750.49 & 3.661 & 1.0064 \\
2.0 & 196 & 685.36 & 3.497 & 1.0059 \\
2.5 & 186 & 545.87 & 2.935 & 1.0059 \\
3.0 & 175 & 474.84 & 2.713 & 1.0059 \\
3.5 & 163 & 452.95 & 2.779 & 1.0073 \\
4.0 & 151 & 428.72 & 2.839 & 1.0068 \\
\bottomrule
\end{tabular}
\end{center}
The residual is concentrated in the low-$\pPar$, near-edge cells; the
high-$\pPar$ tails carry sub-$2\%$ shape disagreement. This is a
robustness check, not a separate method.

\subsubsection{Pull distribution}
\label{sec:pull}

\paragraph{Math.}
For each of the $205$ reported bins,
\begin{equation}
\mathrm{pull}_{i}
\;=\;
\frac{x^{\ours}_{i} - x^{\paper}_{i}}{\sqrt{C^{\tot}_{ii}}},
\qquad
\mathbb{E}[\mathrm{pull}_{i}] = 0,\quad
\mathrm{Var}[\mathrm{pull}_{i}] = 1
\;\text{under the null}.
\end{equation}
Summary diagnostics: pull mean, RMS, the $14\times16$ pull heatmap,
and a pull histogram (in \texttt{compare\_to\_paper\_fullcov.py}).

\paragraph{Worked example (headline numbers).}
With the combined covariance (paper TotalCov $+$ flux-fixed matcorr universe $+$
bootstrap-300 $+$ ML):
\begin{itemize}[nosep]
  \item Pull mean: $\dead{0.069} \to \boxed{0.051}$ (centered on zero; no global offset).
  \item Pull RMS: $\dead{0.466} \to \boxed{0.409}$.
\end{itemize}
The struck values were computed on the pre-flux-fix combined covariance and are
\textbf{retracted}; the live pair is the flux-fixed one recorded in
\texttt{VALIDATION\_LEDGER.md} (\S\ref{sec:combine} provenance note).
Pull RMS $< 1$ means the combined error band is \emph{conservative}
relative to the actual disagreement --- the band is wider than it
needs to be to absorb the observed scatter. No bins have
$|\mathrm{pull}|>3$ in the reported set. Compare with paper-cov-only:
mean $0.089$, RMS $0.598$ --- similar mean (no bias), slightly wider
spread before we include our own uncertainties.

\subsubsection{Shape-only $\chi^{2}$ with Jacobian propagation}
\label{sec:shape}

\paragraph{What it estimates.}
Whether the \emph{shape} of the $(\pT,\pPar)$ cross section matches
the paper, independent of absolute normalization. Useful because the
absolute normalization depends on the flux integral, target nucleon
count, and POT (all of which can drift by $\mathcal{O}(\%)$), whereas
the shape is a pure unfolding diagnostic.

\paragraph{Definition.}
With bin widths $w_{j} = (\Delta\pT)_{j}\,(\Delta\pPar)_{j}$ and total
$T(\mathbf{x}) = \sum_{j\in\text{mask}} w_{j}x_{j}$, the unit-area
shape is
\begin{equation}
s_{i}(\mathbf{x}) \;=\; \frac{x_{i}}{T(\mathbf{x})}.
\end{equation}
By construction $\sum_{i} w_{i} s_{i} = 1$.

\paragraph{Jacobian.}
We need the covariance of $\mathbf{s}$ given the covariance of
$\mathbf{x}$. Differentiating,
\begin{equation}
\frac{\partial s_{i}}{\partial x_{j}}
\;=\;
\frac{\delta_{ij}}{T} - \frac{x_{i}\,w_{j}}{T^{2}}
\;=\;
\frac{1}{T}\big(\delta_{ij} - s_{i}\,w_{j}\big).
\end{equation}
In matrix form (in \texttt{normalize\_xsec\_shape.py}),
\begin{equation}
\mathbf{J} \;=\; \frac{1}{T}\big(\mathbf{I} - \mathbf{s}\,\mathbf{w}\T\big),
\qquad
\mathbf{C}^{\text{shape}} \;=\; \mathbf{J}\,\mathbf{C}^{\paper}\,\mathbf{J}\T.
\end{equation}

\paragraph{Rank reduction.}
$\mathbf{J}$ is rank $n-1$: the unit-area constraint $\sum_{i} w_{i}
s_{i} = 1$ kills one degree of freedom. Algebraically,
$\mathbf{w}\T\mathbf{J} = \mathbf{w}\T(\mathbf{I}-\mathbf{s}\mathbf{w}\T)/T
= (\mathbf{w}\T - \mathbf{w}\T)/T = \mathbf{0}\T$, so the
``all-bins-up'' direction is the kernel of $\mathbf{J}$ and lives in
the null space of $\mathbf{C}^{\text{shape}}$. Accordingly:
\begin{itemize}[nosep]
  \item $\mathbf{C}^{\text{shape}}$ has one zero eigenvalue and rank $n-1$.
  \item $\chi^{2}_{\text{shape}} = (\mathbf{s}^{\ours}-\mathbf{s}^{\paper})\T
        (\mathbf{C}^{\text{shape}})^{+} (\mathbf{s}^{\ours}-\mathbf{s}^{\paper})$
        (pseudo-inverse drops the constrained direction).
  \item $\mathrm{ndf} = n_{\text{bins}} - 1 = 204$ on the
        $205$-bin set.
\end{itemize}

\paragraph{Code pointer.}
\texttt{normalize\_xsec\_shape.py}: \texttt{shape\_jacobian},
\texttt{shape\_chi2}.

\paragraph{Worked example (headline numbers).}
Shape $\chi^{2}/\mathrm{ndf}$ on $205$ bins with the paper's
\texttt{TotalCovariance} (Jacobian-propagated): $\boxed{3.596}$.
This is essentially the absolute-normalization $\chi^{2}/\mathrm{ndf}$
of $3.661$ \emph{minus} the contribution of overall scale. Total-$\sigma$
ratio is $\sigma_{\tot}^{\ours}/\sigma_{\tot}^{\paper} = 1.0111$ ---
$1\%$ scale difference. Shape disagreement at the $\sim 1\%$ level
(median bin ratio $1.0064$, $94\%$ of bins within $10\%$ of the
paper) accounts for the remaining shape residual.

\subsection{Ill-conditioning and rank deficiency}
\label{sec:rank}

\paragraph{The obstruction.}
For publication we would like an \emph{ours-only} $\chi^{2}/\mathrm{ndf}$
--- the data-to-model agreement evaluated with our own covariance
$\mathbf{C}^{\ours} = \mathbf{C}^{\syst} + \mathbf{C}^{\boot} +
\mathbf{C}^{\ML}$ rather than the paper's $V$. That requires inverting
$\mathbf{C}^{\ours}$, and it is rank-deficient:
\begin{equation}
\mathrm{rank}(\mathbf{C}^{\syst}_{\text{sample}}) = 140 / 205,
\qquad
\mathrm{cond}(\mathbf{C}^{\syst}) = 8.7\times10^{8}.
\end{equation}
The deficiency is set by \textbf{band count, not universe count}:
under the MAT formula Eq.~\eqref{eq:matcov} each $\pm1\sigma$ pair is
rank $1$ and the $100$-universe PPFX band is rank $\le 99$ (and
\emph{effectively} rank-$1$: a single normalization mode holds
\SI{99.6}{\percent} of its variance, \S\ref{sec:uq-universes}), so the
$44$-band sum cannot fill the $205$-dimensional bin space. \emph{This is
not a missing-bands problem.} The paper's own systematic-only covariance
($\texttt{TotalCov}-\texttt{StatOnlyCov}$) is rank $120$ --- \emph{lower}
than our $140$ --- so we propagate at least as many independent systematic
directions as the published analysis (we use the same
\texttt{GetStandardSystematics}). The paper's $V$ is full-rank
($\mathrm{rank}=204$, $\mathrm{cond}=1.5\times10^{12}$) purely because it
\textbf{includes the statistical block} ($\texttt{StatOnlyCov}$, rank
$205$, near-diagonal), not because it sums richer systematics. Adding our
$300$-replica bootstrap does the same for us:
$\mathrm{rank}(\mathbf{C}^{\syst}+\mathbf{C}^{\boot})=201$,
$\mathrm{cond}=3.8\times10^{14}$.

\paragraph{The real obstruction is statistical, not rank.}
Even at rank $201$ a defensible single ours-only $\chi^{2}$ is out of
reach --- but for a \emph{statistical-precision} reason, not a missing-band
one. The bootstrap block is smaller in \emph{global scale} by a factor $2.55$ in
square-root trace ($\sqrt{\Tr\mathbf{C}^{\boot}}=1.817\times10^{-40}$ vs
$\sqrt{\Tr\,\texttt{StatOnlyCov}}=4.639\times10^{-40}$), equivalently about
$6.5$ in trace. The paper's full-rank, near-diagonal \texttt{StatOnlyCov} supplies a
variance floor in every reported eigendirection; because the covariance structures
differ, this should \emph{not} be read as a uniform $6.5$-fold per-mode effect. The naive pseudo-inverse then gives
$\chi^{2}/\mathrm{ndf}=252$ for $\mathbf{C}^{\syst}+\mathbf{C}^{\boot}$,
dominated by poorly-determined directions, against the paper's $3.66$ with
its larger stat floor. The lever toward an ours-only $\chi^{2}$ is therefore
the \emph{statistical} block --- not the hunt for more systematic bands.
The gap is \emph{not} a missing MC-statistical term: our bootstrap already
draws independent Poisson($1$) multipliers on \emph{both} data and MC per
replica (\S\ref{sec:bootstrap}; Practical Guide 2507.09582 \S5.1), so
data-stat and MC-stat contribute jointly. With that excluded, the smaller
stat covariance most likely reflects genuinely higher unbinned-OmniFold
statistical efficiency relative to the published binned D'Agostini unfold
(whose iterative regularization amplifies statistical fluctuations), or a
difference in what the paper's \texttt{StatOnlyCov} encompasses ---
the split-bootstrap test below resolves this.

\paragraph{Resolution (data/MC split bootstrap).}
A $200{+}200$-replica \emph{split} bootstrap Poisson-fluctuates the
data weights only (\texttt{boot\_data}) or the MC weights only (\texttt{boot\_mc})
with independent RNGs, so in expectation
$\mathbf{C}^{\boot}=\mathbf{C}^{\boot}_{\mathrm{data}}+\mathbf{C}^{\boot}_{\mathrm{mc}}$.
Three findings settle the question. (i) \emph{Closure holds:}
$\sqrt{\Tr(\mathbf{C}_{\mathrm{data}}+\mathbf{C}_{\mathrm{mc}})}/
\sqrt{\Tr\mathbf{C}^{\boot}}=1.07$, with the variance split $77\%$ data /
$23\%$ MC. (ii) The \emph{data stream alone} gives
$\sqrt{\Tr\mathbf{C}_{\mathrm{data}}}/\sqrt{\Tr\,\texttt{StatOnly}}=0.356$
($\mathbf{C}^{\boot}/\texttt{StatOnly}=0.392=1/2.55$): the $2.5\times$ deficit is
present in the data-only stat error, so it is \emph{not} an MC-statistics
omission --- it reflects genuinely higher unbinned-OmniFold statistical
efficiency than the binned D'Agostini unfold. (iii) There is \emph{also} a
structure difference: our stat block is correlated (mean
$|{\rm off\text{-}diag}\,\rho|=0.108$, leading-eigenvector overlap with
\texttt{StatOnly} $0.46$) whereas the paper's is near-diagonal ($0.009$); scaling
our block to the paper trace still leaves $\chi^{2}/\mathrm{ndf}=76.7$, so the gap
is shape as well as magnitude. \textbf{Consequence:} the correct ours-only
$\chi^{2}$ uses our \emph{own} $\mathbf{C}^{\ours}$, not the paper's stat floor.
The truncation-dependence below remains (the inversion is still regularizer-set),
so we lead with per-bin pulls and the scan; the ours-only goodness-of-fit
is reported on $\mathbf{C}^{\ours}$
via a truncated-spectral $\chi^{2}$ (retain $\lambda>10^{-10}\lambda_{\max}$), the
same convention carried into the 3D generator comparison (\S\ref{sec:rank} is the
2D anchor; the 3D comparison applies it on all $1431$ reported bins).

\paragraph{Regularization scans (all cutoff-dependent).}
On $\mathbf{C}^{\ours}$ over the $205$ reported bins, no regularizer
gives a stable single number:
\begin{center}
\resizebox{\textwidth}{!}{%
\begin{tabular}{@{}lll@{}}
\toprule
Method & Setting & Result \\
\midrule
Eigenvalue truncation & $\lambda/\lambda_{\max} > \{10^{-3},10^{-4},10^{-5}\}$ &
  rank $\{19,49,82\}$; $\chi^{2}/205 = \{0.13,1.12,3.08\}$ \\
Ridge $\mathbf{C}+\varepsilon\lambda_{\max}\mathbf{I}$ & $\varepsilon=\{10^{-3},10^{-4},10^{-5}\}$ &
  $\chi^{2}/205 = \{0.35,1.38,4.30\}$ \\
Diagonal-target shrinkage & $\lambda=0.05$ &
  rank $178\to204$, cond $1.3\times10^{11}$, $\chi^{2}/205 = 3.29$ \\
\bottomrule
\end{tabular}}
\end{center}
Each method spans roughly two decades of $\chi^{2}/205$ across a
two-decade scan of its knob, with no defensible single cutoff. The
useful conclusion is diagnostic, not a number: the naive inverse is
dominated by tiny, poorly-determined modes, and in the retained
high-variance modes the central-value difference is not large.
We therefore report direct per-bin pulls (\S\ref{sec:pull}) and this scan
rather than a single unqualified ours-only $\chi^{2}$.

\paragraph{Bootstrap/paper-stat double-count.}
The combined headline (\S\ref{sec:chisq}) adds $\mathbf{C}^{\boot}$ on
top of the paper's \texttt{TotalCov}, which already contains the
paper's statistical block --- a double count.
\texttt{-{}-subtract-stat} replaces the paper baseline by
$\texttt{TotalCov}-\texttt{StatOnlyCov}$ before adding our covariances,
driving $\chi^{2}/\mathrm{ndf}$ $\chiCombined\to \chiCombinedSubStat$. That \emph{over}-corrects:
our $\sqrt{\Tr\mathbf{C}^{\boot}}=1.83\times10^{-40}$ is far smaller
than the paper's published stat block ($\sim$ the $2.7\times10^{-39}$
scale of \texttt{TotalCov}), so subtracting the latter without a
comparable replacement removes too much. The headline therefore keeps
the double-counted $\chiCombined$ as the well-defined ``what does OmniFold UQ
add on top of the published $V$'' statement, with the de-double-counted
number on file as a secondary diagnostic.

\paragraph{Rank and conditioning: summary.}
The $140/205$ systematic rank is richer than the
paper's systematic-only $120$ (not a missing-bands problem); the Bashyal
flux\,$\leftrightarrow$\,muon-energy joint block is rederived ($\rho=\rhoFluxMuon$)
and carried as a diagnostic --- adding it leaves the systematic rank
unchanged, its directions already lying within the span of the universe
covariance (\S\ref{sec:uq-universes}); the PPFX band is effectively rank-1
(\S\ref{sec:uq-universes}); and the $2.5\times$ smaller bootstrap is explained by
the data/MC split as genuine unbinned-OmniFold efficiency, not a missing MC term.
One paper-independent refinement remains possible: a fitted
$(\Phi_{\text{norm}},\Phi_{\text{shape}},\text{MuonE}_{\text{scale}})$ posterior
from MAT-MINERvA for a self-derived $\rho$. The second --- precedent for the
truncated-spectral ``ours-only'' $\chi^{2}$ convention --- is settled: the
collaboration confirmed (2026-08-02, verbal) that it uses the same
truncated-spectral pseudo-inverse, so the convention here is collaboration
practice rather than an invention of this analysis. That is a confirmation of
practice, not a citation; no published reference has been identified.

\paragraph{Adopted conventions.}
The publication-grade headline uses the combined covariance for
$\chi^{2}/\mathrm{ndf}=\chiCombined$ with the double-count caveat documented;
the per-band rollup is MAT-conformant and byte-for-byte verified
(\S\ref{sec:minerva}); the $1.4\%$ normalization is included as a
rank-1 band; and the ours-only cov is reported via the truncated-spectral
$\chi^{2}$ on $\mathbf{C}^{\mathrm{ours}}$ (rank $201$) with the
regularization scan disclosed.

\subsection{Anatomy of the paper-covariance tension}
\label{sec:tension}

\paragraph{The paradox.}
The central values agree well: $\sigma_{\tot}^{\ours}/\sigma_{\tot}^{\paper}
= 1.011$, median bin ratio $1.006$, and the \emph{per-bin pull RMS}
against the paper covariance is only $0.598$ (\S\ref{sec:pull}).
A diagonal-only $\chi^{2}$ --- ignoring all off-diagonal correlations ---
is therefore
\begin{equation}
\chi^{2}_{\diag}/\mathrm{ndf}
= \frac{1}{n}\sum_{i}\frac{\Delta_{i}^{2}}{C_{ii}}
= \mathrm{RMS}(\mathrm{pull})^{2} + \mathrm{mean}(\mathrm{pull})^{2}
\approx 0.366,
\end{equation}
yet the full-covariance $\chi^{2}/\mathrm{ndf}=3.661$. The tension is a
$10\times$ inflation that lives \emph{entirely} in the off-diagonal /
small-eigenvalue structure of the published \texttt{TotalCov}
(\texttt{diagnose\_tension.py}). The rest of this section dissects it.

\paragraph{Eigenmode decomposition.}
Diagonalize the reported-block paper covariance,
$\mathbf{C}=\sum_{k}\lambda_{k}\mathbf{u}_{k}\mathbf{u}_{k}\T$
($\lambda_{1}\!\ge\!\cdots\!\ge\!\lambda_{n}$), and project the residual
$\boldsymbol{\Delta}$ onto each eigenvector, $c_{k}=\mathbf{u}_{k}\T\boldsymbol{\Delta}$.
Then
\begin{equation}
\chi^{2}=\boldsymbol{\Delta}\T\mathbf{C}^{+}\boldsymbol{\Delta}
=\sum_{k}\frac{c_{k}^{2}}{\lambda_{k}},
\qquad
\text{``eigen-pull''}_{k}\equiv\frac{c_{k}}{\sqrt{\lambda_{k}}},
\end{equation}
verified to reproduce the \texttt{pinv} $\chi^{2}=750.49$ to $3\times10^{-9}$.
The decomposition shows the tension is \textbf{neither a few pathological
null directions nor the absolute smallest eigenvalues}:
\begin{itemize}[nosep]
  \item The top contributing modes sit at $\lambda_{k}/\lambda_{\max}\sim
        10^{-4}$--$10^{-5}$ --- the lower-middle of the spectrum by
        magnitude, but still $\sim7$--$8$ decades above the cond-floor
        ($\lambda/\lambda_{\max}\sim10^{-12}$) --- with standardized-component
        magnitudes $5$--$6$. These are descriptive coordinates in the published
        covariance basis, not calibrated significances for two independent results.
  \item The ten \emph{smallest}-$\lambda$ modes carry only $3\%$ of the
        $\chi^{2}$; about $90$ modes are needed to reach $90\%$. The
        $\chi^{2}$ is broadly distributed, not localized to the
        ill-conditioned tail.
\end{itemize}

\paragraph{Regularization robustness.}
A \texttt{pinv}-rcond scan confirms the tension is not a numerical artifact:
retaining only the $73$ best-determined directions
($\lambda_{k}/\lambda_{\max}>10^{-4}$) still gives $\chi^{2}/\mathrm{ndf}=1.42$;
$139$ directions ($>10^{-6}$) give $2.79$. A rank-truncation scan keeping the
$r$ largest-$\lambda$ modes rises smoothly $0.69\,(r{=}50)\to2.35\,(100)\to
3.30\,(180)\to3.66\,(205)$ --- no late jump, i.e.\ the tension does not ride on
the last few least-determined modes.

\paragraph{It is shape, not normalization.}
Splitting $\boldsymbol{\Delta}$ into a component along the measured
spectrum $\mathbf{x}^{\paper}$ (an overall scale) and the orthogonal
remainder, the scale piece contributes $\chi^{2}_{\mathrm{norm}}=0.10$ and the
shape piece $\chi^{2}_{\text{shape}}=750.1$. The $+1.1\%$ normalization
offset is statistically irrelevant to the tension; it is a correlated
\emph{shape} difference. (Consistent with the Jacobian shape-only
$\chi^{2}/\mathrm{ndf}=3.60$, \S\ref{sec:shape}.)

\paragraph{Where it lives.}
Mapping the dominant eigenvectors and the per-bin contribution
$\Delta_{i}(\mathbf{C}^{+}\boldsymbol{\Delta})_{i}$ back onto the
$14\times16$ grid (\texttt{tension\_mode\_maps.png},
\texttt{tension\_chi2\_map.png}) localizes the tension to the
\textbf{low-$\pT$ cross-section peak ridge} ($\pT\lesssim0.4$ GeV/c,
$\pT$-bins 2/7/10 carry $16/11/12\%$) and the low-$\pPar$ near-edge
cells ($\pPar$-bins 0/1/8). The carrier modes are sign-alternating
\emph{shape} oscillations across the peak, not vertical $\pT$-columns
(which a flux $1/\Phi(\pT)$ normalization error would produce) and not
the $\theta_{\mu}<20^{\circ}$ acceptance boundary. This is exactly where
the flux and Muon\_Energy bands --- the two largest, and the only named
cross-band correlation (Bashyal, \S\ref{sec:rank}) --- dominate the
covariance: tight bin-to-bin correlation there leaves little freedom to
absorb a coherent $\sim1$--$2\%$ shape difference, so a small
OmniFold-vs-IBU shape difference becomes a large $\chi^{2}$.%
\footnote{The pipeline applies one global $\mathrm{POT}_{\mathrm{data}}/\mathrm{POT}_{\mathrm{MC}}$
ratio rather than a per-playlist ratio (KNOWN\_ISSUES J36), and the per-playlist ratios span
$0.1707$--$0.2371$. Because this statement asserts sensitivity at the $1$--$2\%$ level, the
mixture error was measured directly against it: comparing $R_{\mathrm{glob}}\sum_p \mathrm{MC}_p$
with $\sum_p R_p\,\mathrm{MC}_p$ over the twelve playlists changes the reco shape by at most
$0.07\%$ in $\pT$ and $0.14\%$ in $\pPar$ (and $0.03\%$ across the low-$\pT$ ridge), with the
residual $+0.12\%$ appearing as normalization rather than shape. The twelve playlists are nearly
shape-identical, so a large reweighting of them is almost pure normalization. The statement above
is therefore unaffected at the $14$--$30\times$ level.}

\paragraph{A methodological component (estimator and iteration).}
Part of the tension is regularization-dependent --- the same data, same
normalization ($\Sigma\sigma$ within $0.1\%$), unfolded with a different
GBDT backend moves the paper-cov $\chi^{2}$ by $\sim1$ unit while central
values barely change (iteration count, by contrast, is negligible):
\begin{center}
\begin{tabular}{@{}lcc@{}}
\toprule
OmniFold configuration & paper-cov $\chi^{2}/\mathrm{ndf}$ & pull RMS \\
\midrule
exact-GBT, 5 iter (frozen production) & $3.66$ & $0.598$ \\
HistGBT, 5 iter (10-seed mean) & $2.70\pm0.04$ & $0.56$ \\
LightGBM, 5 iter & $2.65$ & $0.563$ \\
LightGBM, 10 iter & $2.66$ & $0.565$ \\
\bottomrule
\end{tabular}
\end{center}
Two things separate cleanly. (i) \textbf{Iteration count is negligible}:
LightGBM $5\to10$ iters moves the $\chi^{2}$ by $+0.01$ ($2.645\to2.657$),
so OmniFold is well-converged at $5$ iters and the tension is \emph{not}
an under-iteration artifact. (ii) \textbf{The GBDT estimator carries a
$\sim1$-unit band}: the frozen exact-GBT production ($3.66$) is the
outlier, while both histogram-binned backends agree at $2.65$--$2.70$
(HistGBT is a $10$-seed mean, range $2.63$--$2.74$; the $\pm0.04$ ML-seed
jitter is $\sim25\times$ smaller than the $0.96$ exact$\to$hist gap, so
that gap is genuinely the estimator, not seed noise). The exact-gradient
trees fit finer shape structure (less implicit smoothing) and so deviate
more from the paper's IBU-regularized result. This is direct evidence
that $\sim1$ $\chi^{2}$-unit of the tension is OmniFold's implicit
regularization rather than a physical data--MC disagreement
(\texttt{tension\_iter/}).

\paragraph{Conclusion.}
The tension is (i) a genuine, broadly-distributed, \emph{correlated shape}
difference --- not a normalization offset and not a small-$\lambda$
inversion artifact --- concentrated in the low-$\pT$ peak where the
flux/Muon\_Energy bands dominate, plus (ii) a $\sim1$-$\chi^{2}$-unit
methodological band from the GBDT estimator/iteration regularization.
Both point to the same structure: the Bashyal
flux\,$\leftrightarrow$\,Muon\_Energy joint block (rederived at
$\rho=\rhoFluxMuon$, \S\ref{sec:rank}) is exactly the covariance structure
governing the region that carries the tension.

\subsection{Pull diagnostics and the open 2D coverage test}
\label{sec:coverage}

\paragraph{What it estimates.}
The existing closure-toy ensemble supports a same-ensemble standardized-pull
calculation that checks whether the toy distribution is approximately
Gaussian. It does not support a valid 2D coverage claim: the MC bootstrap also
fluctuates the stored \texttt{hTruthXSec2D}, so the all-toy truth mean is not an
independent reference. A future disjoint calibration/evaluation test must use
an independently stored unfluctuated truth (or a redesigned toy ensemble).

\paragraph{Procedure (\texttt{sbatch\_coverage\_toys\_MEFHC\_200.sh}).}
Generate $N_{\text{toy}}=200$ closure unfolds, each one combining two
existing modes of \texttt{unfold\_2d\_omnifold\_unbinned.py}:
\begin{itemize}[nosep]
  \item \texttt{-{}-closure}: pseudo-data are MC reco events
        (weighted); a truth marginal is written to
        \texttt{hTruthXSec2D} in every output file.
  \item \texttt{-{}-bootstrap-seed $k$} (seeds chosen to avoid
        collision with the production replicas): per-event Poisson(1)
        draws, same as the production bootstrap.
\end{itemize}
Toy seeds are independent and bins within a toy are correlated. Because the
MC bootstrap multiplies \texttt{w\_truth}, the stored truth marginal also
varies from toy to toy and cannot serve as the independent coverage target.

\paragraph{Math.}
Stack the toys into an $N_{\text{toy}}\times14\times16$ array
$\mathbf{U}^{(t)}$. Here $\bar X^{\text{truth}}$ is the all-toy mean used only
to center this same-ensemble diagnostic. For each reported bin (where
$\bar X^{\text{truth}}_{i}>0$, $i$ indexing the $205$ reported cells):
\begin{align}
\sigma_{i}
\;&=\;
\mathrm{std}\big(\mathbf{U}^{(t)}_{i}; \,\text{ddof}=1\big)
\quad\text{(spread across toys)},
\\
r^{(t)}_{i}
\;&=\;
\frac{U^{(t)}_{i} - \bar X^{\text{truth}}_{i}}{\sigma_{i}}
\quad\text{(per-toy, per-bin standardized residual)},
\\
c_{i}
\;&=\;
\frac{1}{N_{\text{toy}}}
\sum_{t=1}^{N_{\text{toy}}} \mathbf{1}\!\big\{|r^{(t)}_{i}| \le 1\big\}
\quad\text{(same-ensemble per-bin fraction)}.
\end{align}
The reference value of $c_{i}$ for Gaussian standardized residuals
is the central-$\sigma$ probability of a unit normal:
\begin{equation}
c^{\text{expected}}
\;=\;
\Phi(1) - \Phi(-1)
\;=\;
\mathrm{erf}(1/\sqrt 2)
\;\approx\;
0.6827.
\end{equation}
A second diagnostic is the mean absolute pull,
$\langle |r| \rangle$, which for a unit normal is
$\sqrt{2/\pi}\approx 0.798$.

\paragraph{Code pointer.}
\begin{sloppypar}
\texttt{uq/coverage\_toys.py} (rollup); driver
\texttt{sbatch\_coverage\_toys\_MEFHC\_200.sh}; outputs
\texttt{uq/coverage/2d\_xsec\_MEFHC\_5iter\_lgbm\_coverage\_toy*.root}
and \texttt{uq/coverage/coverage\_summary.txt}.
\end{sloppypar}

\paragraph{Worked example (same-ensemble pull diagnostic).}
With $N_{\text{toy}}=200$ on the $205$ reported bins:
\begin{itemize}[nosep]
  \item Fraction with $|r|\leq1$: $\boxed{68.71\%}$ (Gaussian reference $68.27\%$;
        $\Delta = +0.44\%$); median $68.50\%$.
  \item $\langle |r| \rangle = 0.794$ (target $0.798$; within $0.5\%$).
  \item Signed residual mean: $+0.006 \pm 0.082\,\sigma$ (no bias).
  \item Fraction of bins whose within-bin fraction exceeds $65\%$: $\boxed{97.6\%}$ ---
        only $5$ of $205$ bins fall below.
\end{itemize}
The agreement of these same-ensemble quantities with their Gaussian references
shows that the standardized toy residuals are approximately normal.  Because
$\sigma_i$ is estimated from the same toys being scored, this is not a
frequentist coverage measurement and does not certify the bootstrap band.
The attempted 100/100 split-sample reinterpretation is also withdrawn because
it used the all-toy mean of the fluctuated truth histograms; no 2D containment
number is quoted until an independent unfluctuated reference is supplied.

\subsection{Closure tests: detecting bias}
\label{sec:closure}

The variance methods of \S\ref{sec:variance} answer the question
``how much does the unfold wobble.'' Closure tests answer a
different question: ``is the unfold biased --- does it return a value
that is consistently offset from the right answer when we inject a
known truth?'' Bias does not show up in the variance procedures; you
have to look for it.

The three closure tests inject a known deformation into the MC,
unfold, and check that the answer reproduces the deformation. Each
test has its own ``acceptance threshold,'' set comfortably above the
ML-noise floor of $0.166\%$.

\subsubsection{Truth-reweight closure}
\label{sec:closure-tr}

\paragraph{Procedure.}
For a known shape $f(\pT,\pPar)$, multiply both the MC reco weights
\emph{and} the MC truth weights by $f$ at the event level. Unfold the
reweighted pseudo-data against the unchanged (CV) response. The
reference is the reweighted truth marginal --- the unfolded answer
should match it.

\paragraph{Two shapes tested.}
\begin{itemize}[nosep]
  \item \texttt{gauss\_pt}: $f = 1 + A\exp\!\big[-((\pT - p_{T}^{0})/\sigma)^{2}\big]$
        with $A=0.2$, $p_{T}^{0}=0.4$ GeV/c, $\sigma=0.1$ GeV/c.
  \item \texttt{tilt\_pz}: $f = 1 + \alpha(\pPar - \pPar^{\text{ref}})$
        with $\alpha=0.1$, $\pPar^{\text{ref}}=5$ GeV/c.
\end{itemize}

\paragraph{Acceptance.}
Median per-bin $|\text{residual}|\le 1.5\%$.

\paragraph{Code pointer.}
\texttt{uq/closure/closure\_truth\_reweight.py}.

\paragraph{Worked example (1A lgbm 5-iter).}
\begin{itemize}[nosep]
  \item \texttt{gauss\_pt}: median $|\text{residual}|=\boxed{0.046\%}$ (PASS).
  \item \texttt{tilt\_pz}: median $|\text{residual}|=\boxed{0.013\%}$ (PASS).
\end{itemize}
Both are $\sim 100\times$ tighter than the threshold; OmniFold
recovers the injected truth shape essentially perfectly.

\subsubsection{Hidden-variable closure}
\label{sec:closure-hv}

\paragraph{Procedure.}
This is the subtle one. The OmniFold feature set sees $(\pT,\pPar)$
on reco and truth, but \emph{not} the resolution variable
$\dpT \equiv \pT^{\text{reco}}-\pT^{\text{truth}}$. So we inject a
bias on $\dpT$ only --- a kinematic axis the network cannot see ---
and ask: how much of that bias leaks into the measured $(\pT,\pPar)$
marginal?

The pseudo-data weight bump is
\begin{equation}
w^{\text{reco}}
\;\to\;
w^{\text{reco}}\cdot
\Big(1 + A\exp\!\big[-((\dpT - c)/\sigma_{\dpT})^{2}\big]\Big),
\end{equation}
with $c=+0.1$ GeV/c, $\sigma_{\dpT}=0.05$ GeV/c. Truth weights are
\emph{unchanged}, so the reference is the unaltered CV truth
marginal. The amplitude is scanned: $A\in\{0.05,0.10,0.30\}$.

\paragraph{Leakage coefficient.}
Define the signed residual $\bar r(A)$ as the median of the per-bin
signed residual at amplitude $A$. The leakage coefficient is the slope
of a linear fit:
\begin{equation}
\bar r(A) \;\approx\; \kappa_{\text{leak}} \, A
\quad\Rightarrow\quad
\kappa_{\text{leak}}
\;=\;
\frac{\partial \bar r}{\partial A}.
\end{equation}
$\kappa_{\text{leak}}$ tells you what fraction of an unseen
$\dpT$-shaped bias contaminates the visible answer.

\paragraph{Acceptance.}
Median per-bin $|\text{residual}| \le 3.0\%$ (looser than truth-
reweight because $\dpT$ and $\pT$ are weakly correlated, so the
unfolding can never fully kill the leak).

\paragraph{Code pointer.}
\texttt{uq/closure/closure\_hidden\_var.py}.

\paragraph{Worked example (1A lgbm 5-iter, $205$ reported bins).}
\begin{itemize}[nosep]
  \item Per-amplitude median residuals: $0.57\% / 1.14\% / 3.43\%$
        for $A = 0.05/0.10/0.30$ --- linear in $A$, as expected for
        small perturbations.
  \item Linear fit slope: $\kappa_{\text{leak}} = \boxed{0.113}$.
\end{itemize}
About $11\%$ of a hidden $\dpT$ bias leaks into the measured
$(\pT,\pPar)$ marginal --- under the $3\%$ threshold for the largest
realistic deformation ($A=0.30$). Larger leakage would indicate that
$\dpT$ ought to be added to the feature set; the present number is
small enough to be safely ignored.

\subsubsection{Alternative-model closure}
\label{sec:closure-alt}

\paragraph{Procedure.}
Train the OmniFold response on the CV weights. Build the pseudo-data
from an alternative-universe reweighting (e.g.\ \texttt{MaCCQE:0}
$+1\sigma$, or \texttt{Flux:50} PPFX index). The reference is the
in-acceptance-extrapolated alt truth:
\begin{equation}
\tilde X^{\text{ref}}_{i}
\;=\;
X^{\CV,\text{truth}}_{i}\cdot
\frac{X^{\text{alt},\text{in-acc}}_{i}}{X^{\CV,\text{in-acc}}_{i}}.
\end{equation}
This scales the CV truth marginal by the ratio of in-acceptance
counts between alt and CV, isolating the unfolding's response to
a model-shape change from acceptance/efficiency effects. MaCCQE
probes an interaction-model shape change; Flux probes flux
reweighting --- two distinct kinds of model dependence.

\paragraph{Acceptance.}
Median per-bin $|\text{residual}| \le 2.0\%$.

\paragraph{Code pointer.}
\texttt{uq/closure/closure\_alt\_model.py}.

\paragraph{Worked example (1A lgbm 5-iter).}
\begin{itemize}[nosep]
  \item \texttt{MaCCQE:0} ($+1\sigma$): median $|\text{residual}|=\boxed{0.068\%}$ (PASS);
        in-acceptance bias $-1.68\%$ folded out by the extrapolation.
  \item \texttt{Flux:50} (PPFX idx $50$): median $|\text{residual}|=
        \boxed{0.376\%}$ (PASS); in-acceptance bias $-7.32\%$.
\end{itemize}
Both pass with $\gg 5\times$ headroom. The unfolding is robust to
model variations on the order of $\pm 1\sigma$ of GENIE and full PPFX
flux universes.

\subsection{How the pieces fit together}
\label{sec:assemble}

The combined result, in one sentence:
\emph{across $205$ paper-reported bins, the
combined OmniFold $+$ paper error budget gives $\chi^{2}/\mathrm{ndf} =
\chiCombined$ and pull mean/RMS $0.051/0.409$. Because the paper and OmniFold
results share data and their cross-covariance is unavailable, this is an
indicative distance rather than a calibrated goodness-of-fit. The
same-ensemble toy pulls are approximately Gaussian, and the $0.166\%$
ML-noise floor is more than an order of magnitude below the systematic
envelope.}

\subsubsection*{Audit trail of the headline}

\begin{center}
\begin{tabular}{@{}lll@{}}
\toprule
Number & Section & Procedure \\
\midrule
$\chi^{2}/\mathrm{ndf} = \chiCombined$ & \S\ref{sec:chisq} &
  Eq.~\eqref{eq:chi2} with $\mathbf{C}^{\paper}+\mathbf{C}^{\syst}+\mathbf{C}^{\stat}+\mathbf{C}^{\ML}$ \\
$\mathbf{C}^{\syst}$ & \S\ref{sec:universes} &
  $187$ universes, MAT-conformant Eq.~\eqref{eq:matcov} $+$ norm \\
$\mathbf{C}^{\stat}$ & \S\ref{sec:bootstrap} &
  Poisson bootstrap, $300$ replicas, Eq.~\eqref{eq:bootcov} \\
$\mathbf{C}^{\ML}$ & \S\ref{sec:seedscan} & lgbm training-seed variation, $n=10$ trials \\
$\mathbf{C}^{\paper}$ & \S\ref{sec:minerva} & Paper's published \texttt{TotalCovariance} \\
Pull RMS $= 0.409$ & \S\ref{sec:pull} &
  Per-bin standardized residual on combined $\mathbf{C}$ \\
Combined median rel.\ $= 6.865\%$ & \S\ref{sec:combine} &
  $\sqrt{\diag\mathbf{C}^{\tot}}/\mathbf{X}^{\CV}$ \\
Same-ensemble $|r|\leq1$: $68.71\%$ & \S\ref{sec:coverage} &
  Pull diagnostic, $N_{\text{toy}}=200$ \\
\bottomrule
\end{tabular}
\end{center}

\subsubsection*{Standing assumptions}

\begin{enumerate}[nosep]
  \item \textbf{Gaussian approximation.} Pulls and $\chi^{2}$ treat
        each bin as Gaussian. Justified because per-bin entries are
        weighted sums over $\mathcal{O}(10^{4})$ MC events; central
        limit theorem applies.
  \item \textbf{Independence of error sources.}
        Eq.~\eqref{eq:blocksum} requires no cross-correlations between
        ML, statistical, and systematic. The bootstrap pins the GBDT
        seed (so ML noise cancels at first order); universes pin the
        seed too; the training-seed variation pins the bootstrap. Cross-terms are
        higher-order. (The available $N=300$ bootstrap block carries a
        small residual ML leak, \S\ref{sec:bootstrap}.)
  \item \textbf{Bootstrap diagnostic.} \S\ref{sec:coverage}'s
        same-ensemble $68.71\%$ fraction is a Gaussianity check, not a
        half-percent certification of $\mathbf{C}^{\stat}$.
  \item \textbf{Band independence in $\mathbf{C}^{\syst}$.}
        Different physics sources are taken as uncorrelated; the one
        known exception, the Bashyal flux\,$\leftrightarrow$\,muon-E
        joint block, is rederived and carried as a diagnostic --- adding
        it leaves the systematic rank unchanged (\S\ref{sec:rank}).
  \item \textbf{Combined-$\chi^{2}$ double-count.} The headline shares
        systematics with the paper's $V$ and is the ``what OmniFold UQ
        adds on top of $V$'' statement, not an independent ours-only
        test (\S\ref{sec:rank}).
  \item \textbf{Closure passes are valid for the perturbations
        tested.} Nothing guarantees they pass for unphysical
        deformations far outside the training distribution.
\end{enumerate}


\subsection{Multinomial $\to$ independent Poisson bootstrap limit}
\label{app:poisson}

A non-parametric bootstrap with replacement draws $N$ events from the
original $N$-event sample. Let $n_{i}$ be the number of times event
$i$ is drawn. Then $(n_{1},\dots,n_{N}) \sim \mathrm{Multinomial}(N,
\mathbf{p})$ with $p_{i}=1/N$ uniform, $\sum_{i}n_{i}=N$.

The marginal $n_{i}\sim\mathrm{Binomial}(N,1/N)$, whose mean is $1$
and variance $(N-1)/N \to 1$. For large $N$, the binomial converges to
$\mathrm{Poisson}(1)$ (Poisson limit theorem: many trials, small
success probability, fixed expected count).

The full multinomial has the constraint $\sum n_{i}=N$, which
introduces $\mathcal{O}(1/N)$ correlations between the $n_{i}$. In the
large-$N$ limit those correlations vanish and the $n_{i}$ become
\emph{independent} Poisson(1) variables. This is the Poissonization
of the multinomial: the conditioned $N$-event multinomial is a
Poisson process \emph{conditioned} on hitting $N$ counts, and the
unconditional version is independent Poissons.

In practice, the independent-Poisson form is what we implement (one
\texttt{np.random.default\_rng().poisson(1.0)} draw per event), and
the $\sum n_{i}=N$ constraint is dropped. For $N \gtrsim 10^{6}$
(our regime), the relative error from the dropped constraint is
$\mathcal{O}(1/\sqrt{N})$ on the total event count --- far smaller
than the per-bin Poisson variance.

\subsection{The $\pm 1\sigma$ pair covariance: mean-centered (MAT-conformant)}
\label{app:pair}

For a single universe band with two universes $u=0,1$ at $\pm 1\sigma$,
write the two value vectors as $\mathbf{X}_{0}$ and $\mathbf{X}_{1}$,
with universe mean $\langle\mathbf{X}\rangle = \tfrac12(\mathbf{X}_{0}
+\mathbf{X}_{1})$. The MAT-conformant estimator
Eq.~\eqref{eq:matcov} at $N_{u}=2$ is
\begin{equation}
\mathbf{C}^{B}_{\text{mean}}
\;=\;
\tfrac12\sum_{u=0,1}(\mathbf{X}_{u}-\langle\mathbf{X}\rangle)(\mathbf{X}_{u}-\langle\mathbf{X}\rangle)\T
\;=\;
\tfrac14(\mathbf{X}_{0}-\mathbf{X}_{1})(\mathbf{X}_{0}-\mathbf{X}_{1})\T,
\end{equation}
a rank-1 outer product of the \emph{half-difference}. It is mean-centered,
so any common offset of the pair away from $X_{\CV}$ contributes
nothing.

A naive alternative centers on the CV rather than the pair mean, with
$\boldsymbol{\Delta}_{u} = \mathbf{X}_{u}-\mathbf{X}_{\CV}$ and decomposition
$\boldsymbol{\Delta}_{0}=\mathbf{m}+\mathbf{a}$,
$\boldsymbol{\Delta}_{1}=\mathbf{m}-\mathbf{a}$:
\begin{equation}
\mathbf{C}^{B}_{\text{CV-centered}}
\;=\;
\tfrac{1}{2}\big(\boldsymbol{\Delta}_{0}\boldsymbol{\Delta}_{0}\T
                + \boldsymbol{\Delta}_{1}\boldsymbol{\Delta}_{1}\T\big)
\;=\;
\mathbf{m}\mathbf{m}\T + \mathbf{a}\mathbf{a}\T.
\end{equation}
The two agree \emph{only} for a CV-symmetric pair ($\mathbf{m}=0$),
where both reduce to $\mathbf{a}\mathbf{a}\T$. For an asymmetric or
one-sided pair the CV-centered form adds a spurious rank-1
$\mathbf{m}\mathbf{m}\T$ term --- it treats the common offset as a second
variance direction, which mean-centering correctly collapses to near
zero. We mean-center throughout, matching MAT (\S\ref{sec:minerva}); the
bands where the distinction matters are the asymmetric ones (BeamAngle,
MuonResolution, and the one-sided NC band).

\subsection{The full-event statistical covariance: what it is, and what it is not}
\label{app:cstatlimit}

\subsubsection{Construction, stated as a method choice}

The statistical covariance for the full-event extended-phase-space measurement is built from a
$50$-member replica family. Each member repeats the whole analysis --- background-subtracted target
construction, estimator training, and cross-section extraction --- on an independent
$\mathrm{Poisson}(1)$ bootstrap of the measured leg, so that a member's target is the nominal target
multiplied by that draw (\S\ref{app:poisson}). The covariance is the sample covariance of the $50$
resulting cross sections about their own mean, normalised $1/(N-1)$, on the density scale.

This is a method choice and we state it as one. A $\mathrm{Poisson}(1)$ bootstrap of the measured
sample is the standard non-parametric stand-in for measured-statistics uncertainty, but for a
\emph{learned}, non-linear estimator it is not automatic that it inherits the properties it has for a
linear one. We verified that the family is what the construction claims: the multiplicity distribution
reproduces the $\mathrm{Poisson}(1)$ probability mass function for $k=0\ldots5$ to better than
$0.5\%$, the zero-weight fraction is $e^{-1}$ to $2.4\times10^{-4}$ pooled and within $1.9\sigma$ of
$e^{-1}$ in every one of the $18$ populated longitudinal-momentum columns, and the members' refined
targets agree with the nominal's to $0.068\%$ on the rows both retain.

\subsubsection{The measured spread}

Across most of the reported plane the family is tight and the central value sits inside it. In
$6 < p_{\parallel} < \SI{20}{GeV/c}$ it is neither.

\begin{center}
\begin{tabular}{lr}
\hline
median per-cell replica relative s.d., $6<p_{\parallel}<\SI{20}{GeV/c}$ & $67.1\%$ \\
median per-cell replica relative s.d., elsewhere & $8.5\%$ \\
nominal-to-member ratio of unfolded yield in that band, family mean & $3.5969$ \\
s.d.\ of that ratio across the $50$ members & $1.6091$ \\
range across the $50$ members & $1.0859$--$6.8865$ \\
median per-cell nominal-to-family-mean ratio, whole reported plane & $1.000323$ \\
\hline
\end{tabular}
\end{center}

The last row is the context for the others and should not be dropped when they are quoted: over the
reported plane as a whole the central value and the replica mean agree to $0.03\%$ per cell, so the
disagreement is concentrated rather than global. The affected band carries $15.5\%$ of the reported
cross-section integral. An earlier internal statement of the band ratio used a single member's value,
$\dead{5.0467}$ --- \textbf{struck: a $70$th-percentile draw quoted as if typical}; the
family-typical value is $3.5969$ with the spread above.

\subsubsection{What is established}

The discrepancy is localised to the estimator's training, and the extraction is faithful. Comparing
the per-cell unfolded yield before extraction against the extracted cross section, for the nominal and
for a member: in the control region $p_{\parallel} < \SI{6}{GeV/c}$ the two ratios agree to a median
$0.139\%$; in the affected band the ratio of ratios is $1.0000$; and the deficit is separable in
longitudinal momentum alone, for two distinct quantities --- four separate fits, not one pair repeated:

\begin{center}
\begin{tabular}{lcc}
\hline
quantity fitted & $R^{2}$ from $p_{\parallel}$ alone & $R^{2}$ from $p_{T}$ alone \\
\hline
extracted cross-section ratio & $0.868$ & $0.018$ \\
pushed (pre-extraction) yield ratio & $0.839$ & $0.030$ \\
\hline
\end{tabular}
\end{center}

The two rows are close \emph{because} the extraction is faithful. Two candidate mechanisms are retired
by measurement rather than by argument: the replica targets themselves, and the signal factor applied
to truth counts, which the extraction comparison excludes.

The affected cells are not sparse or peripheral. They form a single contiguous band spanning nearly the
full transverse-momentum range, with median reconstruction acceptance $0.859$ --- the highest on the
grid --- carrying $26.5\%$ of all reconstructed-and-accepted signal, and none of them falls below the
low-acceptance threshold that defines this analysis's model-dependent reporting tier. Whatever the
explanation, it is not that the measurement is starved there. For external context, MINERvA's published
double-differential inclusive measurement reports all $84$ bins of this longitudinal-momentum range at
a median fractional statistical uncertainty of $1.55\%$, its worst reported bin anywhere being
$13.1\%$. The two are not the same object --- data statistics on a standard selection against a
bootstrap over a learned estimator on an extended phase space --- and we do not treat the comparison as
decisive. We state it because a reader is entitled to it.

\subsubsection{The fork as it stood when this section was first written}

The family disagrees with the central value by a factor of order three, coherently, in the
best-accepted region of the plane, and at that point two readings were available and not
distinguished. \textbf{(a) The spread is real:} the estimator is genuinely this sensitive to the
measured sample there, and the covariance is right, so the quoted uncertainties in
$6<p_{\parallel}<\SI{20}{GeV/c}$ really are of order $67\%$. \textbf{(b) The proxy is wrong:}
$\mathrm{Poisson}(1)$ resampling gives $36.8\%$ of measured rows zero weight, and a learned estimator
confronted with that much removed support may respond in a way that is a property of the perturbation,
so the covariance would overstate and the construction would need replacing. Both are consistent with
everything above: a genuinely unstable estimator and an invalid perturbation both act during training,
both concentrate where information is scarcest, and both leave the extraction faithful.

\subsubsection{Amendment: the discriminating contrast will not be run, and what the figure therefore is}

\emph{This supersedes part of the preceding subsection, which is retained as written because it records
what was open at the time.}

A discriminating experiment was designed and authorised: $\mathrm{Exponential}(1)$ and
$\mathrm{Poisson}(1)$ share mean and variance and differ only in $P(X{=}0)$, so substituting the former
holds the perturbation's variance fixed while removing its zero-support. It will not be run. Four
independent reviewers recommended against it unanimously, on the ground that \emph{neither outcome would
change what is published}: an unimportant zero atom means publishing the covariance as it stands, and a
decisive one means replacing the resampling scheme with a zero-free one --- which is choosing the noise
model that yields the smaller uncertainty, and is not a defensible basis for tightening an error bar.
The question was closed on argument, not abandoned for cost.

That argument narrows branch (b). The resampled objects are individual reconstructed events. A
$\mathrm{Poisson}(1)$ multiplicity on a row states that, were the data retaken, that event occurred $k$
times; multiplicity zero states that it did not occur, which is what counting noise does to an event
observed once, so the zero atom is not an artifact of the method. Independent $\mathrm{Poisson}(1)$
multiplicities also reproduce exact Poisson fluctuations in every aggregate of those rows
(\S\ref{app:poisson}), which variance-matched continuous alternatives do not: they match the first two
moments and get the discreteness wrong. The resampling scheme is therefore the sampling distribution
rather than a stand-in that might be the wrong one, and ``the proxy is invalid'' is not available as a
reading of the band.

\medskip\noindent
\textbf{What the band figure is.} The $67\%$ is the spread of a \emph{refit} estimator under
\emph{correct} measured-statistics resampling. Two consequences travel with the number wherever it is
quoted. First, it is the sampling uncertainty \emph{of this estimator} --- a learned, iterative
reweighting refit on every replica --- and not a property of the data alone; it includes the estimator's
own sensitivity to which events are present. Second, it is therefore an \emph{upper bound} on how
poorly the data constrain the cross section in that region: a differently regularised estimator could do
better on the same data, and none could do worse purely for want of information. That bound is on the
statistical component only and says nothing about systematics.

\medskip\noindent
\textbf{Status of this amendment.} It is an argument, not a ratified result. It has survived four
independent reviewers including the one who proposed the instrument it retired, which is the strongest
support an argument of this kind has had here, and that is still not a measurement. Ratification rests
with the estimator's owner and the covariance-construction reviewer; the open item that records the
question remains open, and nothing here licenses promoting the central value, quoting the covariance as
verified, or treating the region as settled.

\subsubsection{Limitations of the covariance object itself}

These hold whichever reading is correct. \emph{The covariance was built by one implementation and is
therefore not independently verified.} The construction scope originally called for two blind builders;
the dual-build design was dropped by decision on 2026-08-14 --- a recorded choice, not an omission or an
unmet requirement --- and one builder was produced accordingly. Its specification clauses are
abort-on-failure assertions inside that same builder, which makes them a self-check, and the regression
against the in-tree recipe has power over arithmetic and ordering only, not over domain or centring. No
claim of independent construction or verification is made for this object and none should be inferred
from the clause count.

Every estimated standard deviation carries a $10.1\%$ fractional uncertainty of its own from the family
size alone, $1/\sqrt{2(N-1)}=1/\sqrt{98}$; the retired design document that specified $N=100$ labelled a
family of this size insufficient, and that label stands in the record as written.

The reporting domain and the invertible domain differ, and quoting one for the other is an error. The
measurement is reported on $262$ cells, those with non-zero reconstruction acceptance. The covariance is
supplied on the same domain, but its invertible sub-block is $257$ cells: five cells enter the reported
set in some members and not others, so their apparent variance is partly the reporting mask switching
rather than the cross section moving, and one of the five is reported in only $24$ of the $50$ members.
Those cells are published and flagged and are used in no inversion or goodness-of-fit. The rank is
$49=N-1$, far below $262$, so any use requiring an inverse requires a declared regularisation which this
analysis does not silently supply (\S\ref{sec:rank}). Finally, the covariance is on the density scale
and is centred on the replica mean rather than on the published central value: integrating it requires
bin areas, and the centring means it does not by itself express the offset tabulated above.

\subsubsection{Recomputing these numbers}

Every figure above is derived from committed artifacts, and the per-cell arrays behind the aggregates
are published with them, so a reader can recompute the aggregates and reach a different conclusion. The
family spread and its per-replica values are in
\texttt{state/probe-oi126-band-Rpush-sigma-20260815}; the training-versus-extraction split, its control
region, and the thresholds fixed before that comparison ran are in
\texttt{state/probe-oi126-push-vs-extraction-20260815} and the predeclaration beside it; the
$\mathrm{Poisson}(1)$ verification of the targets is in
\texttt{state/RECEIPT-20260815-oi126-mechanism-narrowing} and
\texttt{state/probe-oi126-zero-fraction-per-column-20260815}; the band geometry, acceptance, share of
integral and the external comparison are in
\texttt{state/RECEIPT-20260815-cstat-tail-geometry-and-weighting-correction}; and the two per-cell
relative standard deviations tabulated above are fields of
\texttt{state/p5a-nominal-vs-cstat-family-percell-20260815}, recomputable from that receipt's own
shipped per-cell array. That receipt carries two retracted share-of-total fields, marked in place;
neither relative-standard-deviation field is among them and no share-of-total figure from it is used
here. The input file the members were built from is measured in
\texttt{state/gate5-source-npz-verified-20260813}, which should be cited for that purpose in preference
to the members' own recorded input digest.
